\documentclass[aspectratio=169,10pt]{beamer}

\usetheme{FGA}

\usepackage[T1]{fontenc}
\usepackage[utf8]{inputenc}
\usepackage{lmodern}
\usepackage{amsmath}
\usepackage{booktabs}
\usepackage{pgfplots}
\pgfplotsset{compat=1.18}

\title[Looking for Twin Deficits]{Looking for the Twin Deficits in a Dollarized Oil-Exporting Economy}
\subtitle{Fiscal shocks, external balance, and oil revenue}
\author{Freddy Garcia-Alban}
\institute{Research seminar}
\date{June 2026}

\pgfplotsset{
  fgaaxis/.style={
    width=0.96\linewidth,
    height=0.52\linewidth,
    axis line style={FGAInk},
    tick label style={font=\scriptsize, text=FGAGray},
    label style={font=\scriptsize, text=FGAInk},
    title style={font=\small\bfseries, text=FGAInk},
    grid=major,
    grid style={FGABone},
    major grid style={FGABone},
    legend style={
      draw=none,
      fill=none,
      font=\scriptsize,
      text=FGAGray,
      at={(0.02,0.98)},
      anchor=north west
    },
    tick style={FGAInk},
  }
}

\begin{document}

\begin{frame}[plain]
  \titlepage
\end{frame}

\section{Motivation}

\begin{frame}{The empirical problem}
  \begin{columns}[T,onlytextwidth]
    \begin{column}{0.57\textwidth}
      \begin{itemize}
        \item Dollarization changes the transmission of fiscal shocks.
        \item Oil revenues move with fiscal space and external accounts.
        \item The object of interest is whether fiscal deficits weaken the trade balance.
      \end{itemize}
    \end{column}
    \begin{column}{0.37\textwidth}
      \fgametric{FGAAubergine}{Dollarized}{No nominal exchange-rate adjustment}
      \vspace{0.10cm}
      \fgametric{FGATeal}{Oil exporter}{Fiscal and external shocks interact}
      \vspace{0.10cm}
      \fgametric{FGACopper}{Twin deficits}{Trade balance response is the test}
    \end{column}
  \end{columns}

  \fgatakeaway{The identification problem is fiscal policy in an oil-linked external sector.}
\end{frame}

\section{Model}

\begin{frame}{Bayesian VAR with sign restrictions}
  \begin{block}{Reduced form}
    \[
      y_t = A_1 y_{t-1} + \cdots + A_p y_{t-p} + u_t,
      \qquad u_t \sim N(0,\Sigma).
    \]
  \end{block}

  \vspace{0.2cm}

  \begin{columns}[T,onlytextwidth]
    \begin{column}{0.31\textwidth}
      \begin{exampleblock}{Priors}
        Minnesota-style shrinkage for macroeconomic dynamics.
      \end{exampleblock}
    \end{column}
    \begin{column}{0.31\textwidth}
      \begin{exampleblock}{Identification}
        Sign restrictions separate fiscal and oil revenue shocks.
      \end{exampleblock}
    \end{column}
    \begin{column}{0.31\textwidth}
      \begin{alertblock}{Target}
        Trade balance and real exchange-rate responses.
      \end{alertblock}
    \end{column}
  \end{columns}
\end{frame}

\section{Results}

\begin{frame}{Fiscal deficits weaken external balance}
  \begin{columns}[T,onlytextwidth]
    \begin{column}{0.63\textwidth}
      \begin{tikzpicture}
        \begin{axis}[
          fgaaxis,
          xlabel={Horizon},
          ylabel={Response},
          xmin=0, xmax=16,
          ymin=-0.8, ymax=0.35,
          xtick={0,4,8,12,16},
          ytick={-0.75,-0.5,-0.25,0,0.25},
          legend entries={Government consumption shock, Tax shock, Key window},
        ]
          \addplot[FGATeal, very thick, mark=none] coordinates {
            (0,0.00) (1,-0.08) (2,-0.15) (3,-0.22) (4,-0.29)
            (5,-0.34) (6,-0.39) (7,-0.45) (8,-0.50) (9,-0.54)
            (10,-0.58) (11,-0.61) (12,-0.63) (13,-0.64) (14,-0.65)
            (15,-0.65) (16,-0.64)
          };
          \addplot[FGAGray, thick, mark=none, opacity=0.75] coordinates {
            (0,0.00) (1,-0.03) (2,-0.08) (3,-0.12) (4,-0.17)
            (5,-0.20) (6,-0.24) (7,-0.27) (8,-0.30) (9,-0.32)
            (10,-0.34) (11,-0.35) (12,-0.36) (13,-0.36) (14,-0.35)
            (15,-0.34) (16,-0.33)
          };
          \addplot[FGACopper, very thick, mark=*, mark size=1.8pt] coordinates {
            (8,-0.50) (9,-0.54) (10,-0.58) (11,-0.61) (12,-0.63)
          };
          \addplot[FGABone, dashed, mark=none] coordinates {(0,0) (16,0)};
        \end{axis}
      \end{tikzpicture}
    \end{column}
    \begin{column}{0.32\textwidth}
      \begin{itemize}
        \item The trade balance deteriorates after expansionary fiscal shocks.
        \item The response is strongest after the medium horizon.
        \item Oil revenue shocks complicate the interpretation of external adjustment.
      \end{itemize}
    \end{column}
  \end{columns}

  \fgatakeaway{Copper should identify the one result the audience should remember.}
\end{frame}

\section{Implications}

\begin{frame}{What the result means}
  \begin{columns}[T,onlytextwidth]
    \begin{column}{0.32\textwidth}
      \begin{block}{Measurement}
        Model fiscal shocks separately from oil-linked revenue movements.
      \end{block}
    \end{column}
    \begin{column}{0.32\textwidth}
      \begin{exampleblock}{Policy}
        External balance constraints matter even without nominal exchange-rate adjustment.
      \end{exampleblock}
    \end{column}
    \begin{column}{0.32\textwidth}
      \begin{alertblock}{Communication}
        The empirical contribution is not a generic VAR, but the institutional setting.
      \end{alertblock}
    \end{column}
  \end{columns}

  \vspace{0.45cm}
  \begin{center}
    {\Large\fgabrand{Aubergine} for structure. \fgatech{Teal} for estimates. \fgakey{Copper} for the result.}
  \end{center}
\end{frame}

\end{document}
